\documentclass[12pt]{article}
\usepackage{amsmath}
\usepackage{amssymb}

\begin{document}

\title{CSE400 - Fundamentals of Probability in Computing\\
Lecture 20: Linear Transformations and Expectations of Multiple RVs}
\author{Dhaval Patel, PhD}
\date{March 19, 2025}
\maketitle

\section*{Outline}

\begin{itemize}
    \item Statistical Representation of Random Vectors
    \item Mean Vector, Correlation, and Covariance Matrices
    \item Linear Transformations
    \item Mapping properties under $Y = AX + b$
    \item Multivariate Gaussian Distributions
    \item From 1D/2D cases to N-dimensional Joint PDFs
    \item Special Cases and Properties
    \item Implications of zero correlation and matrix factorization
\end{itemize}

\section*{Case Study: “Smart City Traffic + Weather System”}

Imagine you are designing an intelligent system in a city like Ahmedabad that predicts traffic congestion using both transportation and climate variables.

Easing ‘pain in the neck’ \\
SVNIT to study 30 junctions in the city with highest traffic congestion, suggest solutions

\subsection*{Step 1: Define N Random Variables}

Let’s define a vector of random variables:

This is your N-dimensional random vector (here $N = 5$).

\subsection*{Step 2: Mean Vector}

\subsection*{Step 3: Covariance Matrix}

\textbf{Interpretation (Intuition and By Definition)}

\begin{itemize}
    \item More rain $\Rightarrow$ more congestion (traffic increases)
    \item More rain $\Rightarrow$ lower speed
    \item More traffic $\Rightarrow$ higher pollution
    \item Higher temperature $\Rightarrow$ pollution disperses
\end{itemize}

covariance matrix = “dependency map of the city”

\subsection*{Step 4: Joint Gaussian Distribution}

We assume:

This means:

The entire system (traffic + climate) behaves in a structured probabilistic way defined by:

\begin{itemize}
    \item Mean (typical conditions)
    \item Covariance (interdependencies)
\end{itemize}

Why Gaussian?

\subsection*{Step 6: Transformation}

Now suppose the Ahmedabad city defines a Congestion Index:

This is a linear transformation:

Inference \# 1

\subsection*{Step 6: Transformation (continued)}

Now suppose the Ahmedabad city defines a Congestion Index:

This is a linear transformation:

If rainfall increases suddenly, what happens?

\begin{itemize}
    \item Speed $\downarrow$
    \item Traffic $\uparrow$
    \item Pollution $\uparrow$
\end{itemize}

Inference \# 2

\subsection*{Inference \# 2 (repeated emphasis)}

If rainfall increases suddenly, what happens?

\begin{itemize}
    \item Speed $\downarrow$
    \item Traffic $\uparrow$
    \item Pollution $\uparrow$
\end{itemize}

\end{document}